\documentclass[11pt]{article}
\usepackage[T1]{fontenc}
\usepackage[a4paper,margin=1in]{geometry}
\usepackage{amsmath,amssymb,amsthm}
\usepackage[hidelinks]{hyperref}
\usepackage{microtype}

\newtheorem*{question}{Source Question 5.1}
\newtheorem*{answer}{Later answer}

\setlength{\parindent}{0pt}
\setlength{\parskip}{0.6em}

\title{Literature Status: Near Unconditionality versus\\
Truncation Quasi-Greediness}
\author{Run \texttt{fa\_banach\_001} --- lane 3}
\date{21 July 2026}

\begin{document}
\maketitle

\begin{center}
\fbox{\parbox{0.9\textwidth}{
\textbf{Status: literature already answered.}
Question~5.1 of arXiv:2209.03445 asks whether a nearly unconditional
basis can fail to be truncation quasi-greedy.  The same authors explicitly
answer this in the affirmative in the separate later paper
arXiv:2305.12253: their Theorems~1.5 and~1.6 produce such a Schauder basis
in a Banach space.
}}
\end{center}

\section{Original question}

F. Albiac, J. L. Ansorena, and M. Berasategui,
\emph{Elton's near unconditionality of bases as a threshold-free form of
greediness}, arXiv:2209.03445, later published in
\emph{Journal of Functional Analysis} \textbf{285} (2023), 110060,
poses the following on arXiv PDF page~19~\cite{source}.

\begin{question}
Is there a nearly unconditional basis that is not truncation quasi-greedy?
\end{question}

The source explains that near unconditionality is equivalent to
quasi-greediness for largest coefficients (QGLC), whereas truncation
quasi-greediness was known to imply QGLC.  The question asks whether the
converse can fail.

\section{Separate later answer}

The same three authors subsequently wrote
\emph{Linear versus nonlinear forms of partial unconditionality of bases},
arXiv:2305.12253, published in \emph{Journal of Functional Analysis}
\textbf{287} (2024), 110594~\cite{answerpaper}.  On arXiv PDF page~5 they
write, immediately before Theorem~1.6, that they solve the question in the
negative---that is, not every QGLC basis is truncation quasi-greedy.

Their Theorem~1.5 states that, for a basis of a quasi-Banach space, the
following are equivalent:
\[
 \text{bounded-oscillation unconditionality}
 \quad\Longleftrightarrow\quad
 \text{truncation quasi-greediness}.
\]
Their Theorem~1.6 then states:

\begin{answer}
There is a Schauder basis of a Banach space that is nearly unconditional
but is not bounded-oscillation unconditional.  Moreover, for every
$1\leq p<\infty$ the example may be chosen with signed fundamental
function comparable to $m^{1/p}$.
\end{answer}

Combining the two theorems, this basis is nearly unconditional but not
truncation quasi-greedy.  It is therefore exactly the counterexample
requested by Question~5.1.  This implication is also the authors' stated
purpose in the paragraph introducing Theorem~1.6, so the identification is
explicit rather than an agent-inferred consequence.

\section{Scope and search status}

The later construction completely answers the existence question in
Question~5.1, and does so already in a Banach space.  It does not by itself
settle Question~5.2 of the source concerning equivalent renormings and
uniform QGLC constants.

A bounded search through 21 July 2026 checked the run registry, local arXiv
corpus, exact question wording, QGLC/truncation terminology, and later work
by the source authors.  The decisive match is arXiv:2305.12253 and its
published JFA version.  No new mathematical result is claimed in this
packet.

\begin{thebibliography}{9}
\bibitem{source}
F. Albiac, J. L. Ansorena, and M. Berasategui,
\emph{Elton's near unconditionality of bases as a threshold-free form of
greediness}, J. Funct. Anal. \textbf{285} (2023), 110060;
arXiv:2209.03445.

\bibitem{answerpaper}
F. Albiac, J. L. Ansorena, and M. Berasategui,
\emph{Linear versus nonlinear forms of partial unconditionality of bases},
J. Funct. Anal. \textbf{287} (2024), 110594;
arXiv:2305.12253.
\end{thebibliography}

\end{document}
